%% kapa: exam=1 round=37 year=2026 period=1 form=A subject=economics subject_ko=경제학원론 seq=16 number=56
%% source: Q-Net 감정평가사 기출문제 (cst00309) artlSeq=5258924 file=raw/1008/5258924/2026년 제37회 감정평가사 1차 1교시.pdf
%% conversion: kordoc 4.24.51 -> tools/gen1.py (auto); review=proofread-sonnet; flags=none
\begin{question}{56}
사회후생에 관한 설명으로 옳지 않은 것을 모두 고른 것은?
\begin{framedtext}
ㄱ. 차선의 이론은 부분적이고 점진적인 해결책이 최적임을 보여준다. \par
ㄴ. 롤즈(J. Rawls)적 가치판단을 반영한 사회무차별곡선은 L자 모양이다. \par
ㄷ. 공리주의적 가치판단을 반영한 사회무차별곡선은 곡선의 형태이다. \par
ㄹ. 애로우(K. Arrow)의 불가능성 정리에서 파레토원칙은 과반수제를 의미한다. \par
\end{framedtext}
\begin{choices}
  \item ㄱ, ㄴ
  \item ㄴ, ㄷ
  \item ㄱ, ㄴ, ㄷ
  \item ㄱ, ㄷ, ㄹ
  \item ㄴ, ㄷ, ㄹ
\end{choices}
\end{question}
\answer{4}
